\section*{ЗАКЛЮЧЕНИЕ}
\addcontentsline{toc}{section}{ЗАКЛЮЧЕНИЕ}

В ходе выполнения работы был разработан корпоративный мессенджер, полностью независимый от зарубежных поставщиков и санкционных ограничений. Проведённое исследование подтвердило важность автономной системы обмена сообщениями для российских компаний, обеспечивающей стабильную работу внутренних коммуникаций и защиту данных.

Основные результаты работы:

\begin{enumerate}
	\item В ходе анализа предметной области были рассмотрены требования к корпоративному мессенджеру, выявлены риски использования зарубежных платформ и определены ключевые принципы автономности системы.
	\item Разработаны концептуальная, логическая и физическая модели данных корпоративного мессенджера, обеспечивающие точное отражение бизнес-процессов обмена сообщениями.
	\item Спроектирована реляционная схема базы данных с учётом особенностей управления пользователями, передачи сообщений и хранения вложений, что позволило создать эффективную систему обработки данных.
	\item Реализована модель данных в СУБД SQLite, протестирована работа системы, проведено комплексное тестирование, включая модульное и интеграционное тестирование, что подтвердило её корректность и устойчивость.
\end{enumerate}

Разработанная система представляет собой комплексное решение для корпоративного обмена сообщениями, ориентированное на автономную работу и защиту данных. Были реализованы механизмы отправки сообщений, управление группами пользователей, проверка статуса сообщений и обработка вложений. Благодаря интеграции всех компонентов обеспечивается единое информационное пространство, способствующее эффективному обмену данными внутри компании.

Все требования, объявленные в техническом задании, были полностью реализованы, а задачи, поставленные в начале разработки проекта, решены в полном объёме. 

Перспективой дальнейшей разработки является расширение функционала мессенджера, добавление поддержки масштабируемых баз данных и внедрение дополнительных механизмов безопасности.
